\documentclass[10pt,a4paper]{amsart}
\usepackage[margin=19mm]{geometry}
\usepackage{amsmath,amssymb,mathtools}
\usepackage[scr=rsfs,cal=euler]{mathalfa}
\usepackage{xfrac}
\usepackage[unicode,hidelinks,bookmarks=false]{hyperref}
\newcommand{\ie}{i.e.\ }
\newcommand{\cf}{cf.\ }
\newcommand{\resp}{respectively\ }
\newcommand{\Subsection}{\subsection}
\providecommand{\nobreakdash}{\nobreak\hbox{-}}
\setlength{\emergencystretch}{2em}
\multlinegap0pt
\usepackage{amssymb}
\usepackage{mathtools}
\newtheorem{theorem}{Theorem}
\title{Hierarchical Solomonoff Induction: An Unbounded Machine Learning Model}
\author{Nathan Young}
\date{Source: arXiv:2608.01005v1 (CC BY 4.0)}
\begin{document}
\maketitle

\noindent\emph{Source and licence.} Hierarchical Solomonoff Induction: An Unbounded Machine Learning Model. Version arXiv:2608.01005v1 (CC BY 4.0), distributed by arXiv under the Creative Commons Attribution 4.0 International licence, http://creativecommons.org/licenses/by/4.0/, as recorded in the arXiv licence metadata for this exact version and shown on the abstract page https://arxiv.org/abs/2608.01005v1. Full licence notice: \texttt{licenses/arxiv-2608.01005-CC-BY-4.0.txt}.\par\smallskip

\noindent\textit{Editorial scope.} Counted unit: the article's theorem thm:enums-of-ums=ums (source0.tex lines 268--273). The article's complete original proof of it is delivered with it (source0.tex lines 274--294, one \texttt{proof} environment of 21 lines). No mathematics is written, repaired or re-derived: the statement and the proof are the article's own lines, in the article's own notation, and no retained line is substituted or re-flowed.  No further source-local context is needed: every cross-reference of every retained line resolves inside the two delivered spans.  Nothing was omitted from any retained span, so the package carries no omission record.  No retained line contains a citation, so the wrapper carries no bibliography and no bibliography entry.  No retained line contains a figure environment, an image-inclusion command or drawing code, so no image is delivered and the package declares no asset dependency.  Editorial formatting only, in addition: an AMS \texttt{amsart} preamble that restates the article's own declarations and macros used by the retained lines, namely the environments \texttt{theorem} on separate counters, together with the standard packages those lines need.  The retained environments print in this document as Theorem 1 in order of appearance, whereas the article prints the same items with the numbering of its own full document.  The complete native source of the article as distributed in the arXiv e-print for this exact version is bundled unchanged as \texttt{sources/206586/source0.tex}.
\begin{theorem}
\label{thm:enums-of-ums=ums}
Let $\{\nu_i\}_{i=1}^\infty$ be any enumeration of enumerable semimeasures in which every universal mixture $\xi\in\mathcal{U}_\xi$ appears.
Let $\mathcal{C}_\nu$ be the class of all mixtures $c\coloneq\sum_{i\in\mathbb{N}} w^{(c)}(i)\nu_i$ whose weight functions $w^{(c)}:\mathbb{N}\to\mathbb{R}$ are lower-semicomputable with $w^{(c)}(i)>0\;\forall i\in\mathbb{N}$ and $\sum_{i\in\mathbb{N}}w^{(c)}(i)\leq1$.
Then $\mathcal{C}_\nu=\mathcal{U}_\xi$.
\end{theorem}

\begin{proof}
We take any $c\in\mathcal{C}_\nu$ and any $\xi\in\mathcal{U}_\xi$, and show that each can be used to construct an instance of their own class equivalent to the other.

We first prove that $\mathcal{C}_\nu\subseteq\mathcal{U}_\xi$.
Since $\nu$ includes all universal mixtures, $\xi$ must appear at some index $i_\xi^{(\nu)}$. Define a new enumeration of semimeasures $\nu'$ that alternates between the outputs of $\nu$ and $\xi$'s own semimeasure enumeration $\nu^{(\xi)}$, skipping $\xi$ in $\nu$. Note that as a universal mixture, $\xi$ appears in its own enumeration.
Give each semimeasure in $\nu$ its corresponding weight in $w^{(c)}$ (again skipping $\xi$), and give each semimeasure in $\nu^{(\xi)}$ its weight in $w^{(\xi)}$ multiplied by $w^{(c)}(i_\xi^{(\nu)})$.
This defines a measure with lower-semicomputable weights on every enumerable semimeasure --- that is, a universal mixture --- equal to $c$. Then $c\in\mathcal{U}_\xi$ and $\mathcal{C}_\nu\subseteq\mathcal{U}_\xi$.

We then prove that $\mathcal{C}_\nu\supseteq\mathcal{U}_\xi$.
Since $\nu^{(\xi)}$ includes all enumerable semimeasures, $c$ must appear at some index $i_c^{(\xi)}$.
Compute $w^{(\xi)}(i_c^{(\xi)})$ from below to find some positive rational $q\leq w^{(\xi)}(i_c^{(\xi)})$ and fix $\delta = \frac{q}{2}$ so that $0<\delta<w^{(\xi)}(i_c^{(\xi)})$.
Define a new universal mixture $\xi'$ with identical semimeasure enumeration to $\xi$ but a new weight function ${w^{(\xi)}}'$:
\[{w^{(\xi)}}'(i) = \begin{cases} \frac{w^{(\xi)}(i)}{1-\delta} \text{ if } i \neq i_c^{(\xi)} & \\ \frac{w^{(\xi)}(i)-\delta}{1-\delta} \text{ if } i = i_c^{(\xi)} \end{cases}\]
As a universal mixture, $\xi'$ must appear in $\nu$ at some index $i_{\xi'}^{(\nu)}$.
We may then define a new weight function ${w^{(c)}}'$ over $\nu$:
\[{w^{(c)}}'(i) = \begin{cases} 1-\delta + \delta w^{(c)}(i) \text{ if } i = i_{\xi'}^{(\nu)} & \\ \delta w^{(c)}(i) \text{ if } i \neq i_{\xi'}^{(\nu)} \end{cases}\]
Multiplying $\xi'$ by $1-\delta$ recovers the original weight function of $\xi$, except for the subtraction of $\delta$ from its weight on $c$; this is recovered by every $\nu_i$ receiving its weight in $w^{(c)}$ multiplied by $\delta$.
This defines a measure with lower-semicomputable weights over $\nu$ --- that is, a member of $\mathcal{C}_\nu$ --- equal to $\xi$. Then $\xi\in\mathcal{C}_\nu$ and $\mathcal{C}_\nu\supseteq\mathcal{U}_\xi$.

We then have $\mathcal{C}_\nu\subseteq\mathcal{U}_\xi\subseteq\mathcal{C}_\nu$, which gives $\mathcal{C}_\nu=\mathcal{U}_\xi$.
\end{proof}

\end{document}
